% !PRINTED 299
\weberpage{299}
%<<CONT
in $[z_1+\alpha_1, z_2+\alpha_2, \ldots, z_k+\alpha_k]$ \"ubergeht. Demnach k\"onnen wir die Permutation $A$ auch so bezeichnen:
\weq{4}{\begin{pmatrix} z_1, & z_2, & \ldots, & z_k \\ z_1+\alpha_1, & z_2+\alpha_2, & \ldots, & z_k+\alpha_k \end{pmatrix},}
und man erh\"alt die ganze Gruppe $Q$, wenn man $\alpha_1, \alpha_2, \ldots, \alpha_k$ je ein volles Restsystem nach dem Modul $p$ durchlaufen l\"asst.

Die Gruppe $Q$, deren Bildungsweise jetzt festgestellt ist, muss ein Normaltheiler von $P$ sein, und daraus l\"asst sich eine allgemeine Form herleiten, in der alle Permutationen von $P$ enthalten sein m\"ussen.

Dazu bed\"urfen wir eines H\"ulfssatzes \"uber die analytische Darstellung der Permutationen, den wir zun\"achst ableiten.

\begin{center}
\S.\,77.\\
\gesperrt{Analytische Darstellung der Permutationen.}
\end{center}

\gesperrt{H\"ulfssatz. Wenn bei irgend einer Permutation der Gr\"ossen (2), \S.\,76, $z_1$ in $z_1'$, $z_2$ in $z_2'$, \ldots, $z_k$ in $z_k'$ \"ubergeht, so kann man immer
\weq{1}{
\begin{aligned}
z_1' &\equiv \varphi_1(z_1, z_2, \ldots, z_k)\\
z_2' &\equiv \varphi_2(z_1, z_2, \ldots, z_k)\\
&\;\cdots\cdots\cdots\cdots\\
z_k' &\equiv \varphi_k(z_1, z_2, \ldots; z_k)
\end{aligned}
\quad (\mathrm{mod}\ p).
}
setzen, worin $\varphi_1, \varphi_2, \ldots, \varphi_k$ ganze rationale Functionen der Variablen $z_1, z_2, \ldots, z_k$ mit ganzzahligen Co\"efficienten bedeuten, die in Bezug auf keine der Variablen den Grad $p - 1$ \"ubersteigen.}

Der Satz ist eine Verallgemeinerung des entsprechenden Satzes im \S.\,180 des ersten Bandes. Dort ist der Satz f\"ur $k = 1$ bewiesen. Dem allgemeinen Beweise, der auf der vollst\"andigen Induction beruht, schicken wir folgende Erw\"agungen voraus.

Da wir in (1) f\"ur die Argumente $z_i$ alle Combinationen von $k$ Reihen nach dem Modul $p$ zu setzen haben, so ergeben sich $p^k$ solcher Systeme. Jede der Functionen $\varphi$ hat $p^k$ unbestimmte Co\"efficienten, und wenn also die $z_k'$ gegeben sind, so erhalten wir $kp^k$ lineare Congruenzen f\"ur ebenso viele unbekannte Zahlen, n\"amlich die Co\"efficienten der $\varphi_k$. Es ist nur noch nachzuweisen,
%CONT>>
